\section{The Benchmark, the Controls and the Probes}
\label{sec:method}

This section describes the benchmark whose metric is under examination and the
two families of submission with which we measured it. Neither family reads a
video frame.

\subsection{The competition}
\label{sec:competition}

Measurements were made on Zero-shot Accident Anticipation, a public competition
hosted on Kaggle~\cite{autopilot2026}. It ran from 17 February to 15 April
2026, drawing 59 entrants of whom 20 submitted, forming 13 teams across 194
submissions. Late submission remains open and is scored against the same
splits, which is how the experiments of Section~\ref{sec:leaderboard} were run
after the closing date. Every figure quoted here was read from the
competition's own pages on 31 August 2026.

Two properties make it a suitable instrument. The organisers hold the labels
and compute every score, so nothing reported here depends on our own
implementation of a metric. And a submission is a plain list of numbers per
clip, so it need not contain a model at all --- which is what makes a
video-blind control submittable.

The corpus is a curated 1{,}417-clip subset of MM-AU~\cite{fang2024mmau}, every
clip exactly 150 frames at 30\,fps. The published test file carries an
identifier, a video identifier, a start and end frame, and a caption. There is
no label, no accident time and no alert time.

One consequence deserves emphasis, because the rest of this paper turns on it.
Average precision, area under the ROC curve, a false-positive rate, and any
time-to-accident measured to a true onset are \emph{not computable by any
entrant}. An entrant wanting to measure progress locally must substitute
something, and the only quantities computable from a submission alone are
properties of the risk curve's own shape --- the frame at which it crosses the
threshold, whether it falls back, how steeply it rises. We did exactly this,
and Section~\ref{sec:leaderboard} measures what that proxy was worth. A second
consequence we identify without quantifying: the 79 distinct captions describe
outcomes (``lead vehicle stops'', 146 clips; ``a vehicle controls loss'', 144),
and they are supplied at inference time, so a method conditioned on them has
access to the outcome it is asked to anticipate.

\subsection{The metric, as published}
\label{sec:metric}

The evaluation page defines the score as a weighted sum of four quantities,
\begin{equation}
\text{score} = w_{\text{AP}}\text{AP} + w_{\text{AUC}}\text{AUC}
 + w_{\text{TTA}}\text{TTA@0.5} + w_{\text{STTA}}\text{STTA@0.5}
\end{equation}
with the weights fixed by the organisers and not disclosed. Positives are
accident clips, negatives ordinary driving. The timing terms are read at a risk
threshold of $0.5$:
\begin{align}
\text{TTA@0.5} &= \max\{\, t_{ai} - t_a \mid p_{t_a} > 0.5,\ 0 \le t_a \le t_{ai} \,\} \\
\text{STTA@0.5} &= \max\{\, t_{ai} - t'_a \mid p_t > 0.5\ \forall\, t \in [t'_a, t_{ai}] \,\}
\end{align}
where $t_{ai}$ is the accident start frame and $t_a$ the first frame at which
the risk score exceeds the threshold. STTA additionally requires the score to
remain above the threshold continuously from $t'_a$ to $t_{ai}$.

Three observations follow from these definitions alone, before any measurement.

\emph{First}, AP and AUC are bounded on the unit interval, while TTA and STTA
are counted in frames and bounded only by $t_{ai}$ --- hence by how the corpus
was windowed. The four terms are not commensurable, and their sum inherits the
scale of the larger pair.

\emph{Second}, both timing terms are maximised by the same trivial submission.
A risk score exceeding $0.5$ at frame~0 and never falling has $t_a = 0$ and
satisfies the continuity condition on every clip, attaining
$\text{TTA} = \text{STTA} = t_{ai}$, the largest value each clip admits. No
submission can exceed it; one can only match it, and then compete on AP and
AUC.

\emph{Third}, neither timing term reads the magnitude of the score, only
whether it stands above $0.5$. A submission that is barely committed and one
that is certain receive identical timing credit, so the metric offers no
incentive to calibrate.

\subsection{Video-blind controls}
\label{sec:controls}

A video-blind submission is one whose risk curve is a function of the frame
index alone: identical for every clip, computed without opening a frame. We use
a family rather than a single curve, because one flattering curve invites the
objection that the comparison was chosen to flatter it. The family comprises
constants at $0.51$ and $0.99$; a curve held at $0.49$, which never crosses; a
linear ramp; logistic curves centred at frames 60, 75, 90 and 105; step
functions rising from $0.49$ to $0.51$ at frames 0, 10, 25, 50, 75, 100, 125
and 140, plus a calibration step at frame 76; a curve that crosses at frame~0
then drops back, and one that crosses and dips across frames 50--59; and two
graded curves. That is 21 distinct shapes submitted as 25 files --- four were
sent twice, and the repeats agree to within $3\times10^{-5}$, which is a check
on the scoring service and not only on our bookkeeping.

Before spending a submission slot we regenerated the organisers' sample file
from its own parsed values and compared it byte for byte against the
distributed file: identical. A scoring anomaly therefore cannot be attributed
to our file writer.

\subsection{Probes that isolate the terms}
\label{sec:probes}

The decomposition rests on one observation. Every video-blind submission is
identical across clips, so --- provided the scorer reduces each curve to one
clip-level score by any deterministic function of that curve --- every member
of the family produces an all-way tie among clips, and therefore earns
identical AP and identical AUC. Whatever those terms are worth, they are worth
the same to all of them, and \emph{they cancel exactly in any difference
between two members}.

That turns the leaderboard into an instrument. Three probes suffice. A curve
held at $0.49$ never crosses, so both timing terms are zero and its score is
the discrimination floor alone. A curve at $0.51$ in frame~0 and $0.49$
thereafter crosses once and then falls, earning the time-to-accident term but
forfeiting the stable one. A constant $0.51$ earns both. Differencing the three
gives each term separately. A fourth probe, holding $0.51$ throughout except
for a dip across frames 50--59, moves the stable term's reference point while
leaving the first crossing untouched, and serves as a consistency check.

The step family measures the score's dependence on the crossing frame directly,
giving a second and independent route to the same quantity. Two routes to one
number is what makes the result safe to report.
